
\documentclass[11pt]{article}

\usepackage[T1]{fontenc}
\usepackage{newtxtext,newtxmath}
\usepackage[margin=1in]{geometry}
\usepackage{microtype}
\usepackage[hidelinks]{hyperref}
\usepackage{setspace}
\usepackage{parskip}
\usepackage{enumitem}
\usepackage{longtable}
\usepackage{booktabs}
\usepackage{array}

\setlength{\parindent}{0pt}
\setstretch{1.10}
\setlist[itemize]{leftmargin=1.5em,itemsep=0.25em,topsep=0.25em}

\begin{document}

\begin{center}
{\Large\bfseries Supplementary Information}

\vspace{0.55cm}

{\large\bfseries
A Multi-Track Simulation of Laser-Based Powder Bed Fusion of Metals
(PBF-LB/M) with Novel Double Conical Heat Source to Predict the
Effects of Laser Turnaround Time on Pad Dimensions and Cooling Rates
for NIST AM-Bench 2025 Challenge
}

\vspace{0.65cm}

Rahul Singha Rathun$^{1}$,
Rakibul Islam Kanak$^{2}$,
Badhon Kumar$^{2}$,
Abdullah Al Amin$^{1,*}$

\vspace{0.45cm}

$^{1}$Department of Mechanical and Aerospace Engineering,
School of Engineering, University of Dayton,
Dayton, OH 45469, USA

\vspace{0.20cm}

$^{2}$Department of Mechanical Engineering,
Bangladesh University of Engineering and Technology,
Dhaka 1000, Bangladesh

\vspace{0.45cm}

$^{*}$\textbf{Corresponding author:} Abdullah Al Amin, Assistant Professor\\
Department of Mechanical and Aerospace Engineering, School of Engineering,
University of Dayton\\
Email: \href{mailto:aamin1@udayton.edu}{aamin1@udayton.edu}

\vspace{0.45cm}

\textit{The International Journal of Advanced Manufacturing Technology}
\end{center}

\vspace{0.6cm}

\section*{Overview}

This Supplementary Information document describes the electronic
supplementary materials accompanying the manuscript. The supplementary
package contains 23 files: seven transient simulation videos, seven
temperature-history videos, eight spreadsheet files containing supporting
numerical data, and one image illustrating the steady-state condition.
In the manuscript and the journal submission system, these materials are
cited consecutively as \textit{Online Resource 1} through
\textit{Online Resource 23}.

The supplementary materials provide additional visualization and numerical
data supporting the model development, calibration, validation, and analysis
presented in the manuscript. Detailed descriptions of the individual online
resources are provided in the following sections.


\section*{S1. Transient Simulation Videos}

Online Resources 1--7 provide transient contour animations of the numerical
simulations presented in the manuscript. Each video includes contours of
temperature, liquid fraction, volumetric laser heat-source intensity,
radiative heat-loss source term, and convective heat-loss source term.
Temperature is reported in K, liquid fraction is dimensionless, and the
volumetric heat-source intensity and heat-loss source terms are reported in
W/m$^3$.

For each field, contours are presented in the $x$--$z$ plane, representing
a longitudinal section along the laser scanning direction, and in the
$x$--$y$ plane, representing the top view of the build surface. The videos
illustrate the temporal evolution of the thermal field, melt-pool formation
and solidification, applied volumetric heat source, and heat losses during
laser scanning. The opening screen of each video provides the corresponding
simulation conditions and numerical parameters, including mesh information,
time-step size, track length, laser power, scan speed, and other relevant
case-specific parameters.


\textbf{Online Resource 1 (ESM\_1.mp4).}\\
Transient thermal simulation of the NIST AM-Bench 2022 single-track
Case 0 using the conventional single-conical volumetric heat-source
model. This simulation is used in the single-track validation presented
in the main manuscript and serves as the reference heat-source model
for comparison with the proposed Double Conical Heat Source (DCHS).
The video presents the transient evolution of the temperature field,
liquid fraction, volumetric laser heat-source intensity, and radiative
and convective heat-loss source terms. The simulation was performed
with a laser power of 285 W, a scan speed of 960 mm/s, and a track
length of 2 mm. Detailed computational and processing parameters are
provided on the opening screen of the video.

\medskip

\textbf{Online Resource 2 (ESM\_2.mp4).}\\
Transient thermal simulation of the NIST AM-Bench 2022 single-track
Case 0 using the proposed Double Conical Heat Source (DCHS) model.
This simulation is used for the single-track validation of the proposed
model in the main manuscript and is directly compared with the
single-conical heat-source simulation provided in Online Resource 1.
The video presents the transient evolution of the temperature field,
liquid fraction, volumetric laser heat-source intensity, and radiative
and convective heat-loss source terms. The simulation uses the same
process conditions as Online Resource 1, with a laser power of 285 W,
a scan speed of 960 mm/s, and a track length of 2 mm. Detailed
computational and processing parameters are provided on the opening
screen of the video.

\medskip

\textbf{Online Resource 3 (ESM\_3.mp4).}\\
Transient thermal simulation of the NIST AM-Bench 2022 XPad
multi-track benchmark using the conventional single-conical
volumetric heat-source model. This simulation is used in the
multi-track XPad validation presented in the main manuscript and
serves as the reference case for comparison with the proposed Double
Conical Heat Source (DCHS) model. The video presents the transient
evolution of the temperature field, liquid fraction, volumetric laser
heat-source intensity, and radiative and convective heat-loss source
terms during the simulation of 10 bidirectional tracks. The simulation
was performed with a laser power of 285 W, a scan speed of 960 mm/s,
a track length of 2 mm, and a hatch spacing of 110~$\mu$m. Laser
turnaround times of 5 ms and 40 ms were applied after odd- and
even-numbered tracks, respectively. Detailed computational and
processing parameters are provided on the opening screen of the video.

\medskip

\textbf{Online Resource 4 (ESM\_4.mp4).}\\
Transient thermal simulation of the NIST AM-Bench 2022 XPad
multi-track benchmark using the proposed Double Conical Heat Source
(DCHS) model. This simulation is used in the multi-track XPad
validation of the proposed model presented in the main manuscript
and is directly compared with the conventional single-conical
heat-source simulation provided in Online Resource 3. The video
presents the transient evolution of the temperature field, liquid
fraction, volumetric laser heat-source intensity, and radiative and
convective heat-loss source terms during the simulation of 10
bidirectional tracks. The simulation was performed with a laser
power of 285 W, a scan speed of 960 mm/s, a track length of 2 mm,
and a hatch spacing of 110~$\mu$m. Laser turnaround times of 5 ms
and 40 ms were applied after odd- and even-numbered tracks,
respectively. Detailed computational and processing parameters are
provided on the opening screen of the video.

\medskip

\textbf{Online Resource 5 (ESM\_5.mp4).}\\
Transient thermal simulation of the NIST AM-Bench 2025 single-track
case using the proposed Double Conical Heat Source (DCHS) model.
This single-track case was provided as part of the AM-Bench 2025
competition and was used to calibrate the DCHS model prior to its
application to the multi-track Challenge 7 cases presented in the
main manuscript. The video presents the transient evolution of the
temperature field, liquid fraction, volumetric laser heat-source
intensity, and radiative and convective heat-loss source terms during
laser scanning. The simulation was performed with a laser power of
285 W, a scan speed of 960 mm/s, a track length of 2 mm, and a laser
spot diameter of 72~$\mu$m. Detailed computational and processing
parameters are provided on the opening screen of the video.

\medskip

\textbf{Online Resource 6 (ESM\_6.mp4).}\\
Transient thermal simulation of the NIST AM-Bench 2025 Challenge 7
multi-track case with a laser turnaround time of 5 ms using the proposed
Double Conical Heat Source (DCHS) model. This case constitutes one of
the primary multi-track prediction cases of the AM-Bench 2025 competition
and was simulated using the DCHS model calibrated from the corresponding
single-track case provided in Online Resource 5. The video presents the
transient evolution of the temperature field, liquid fraction, volumetric
laser heat-source intensity, and radiative and convective heat-loss source
terms during 15 bidirectional tracks. The simulation was performed with
a laser power of 285 W, a scan speed of 960 mm/s, a track length of 1 mm,
a hatch spacing of 110~$\mu$m, and a laser spot diameter of 72~$\mu$m.
A turnaround time of 5 ms was applied between consecutive tracks.
Detailed computational and processing parameters are provided on the
opening screen of the video.

\medskip

\textbf{Online Resource 7 (ESM\_7.mp4).}\\
Transient thermal simulation of the NIST AM-Bench 2025 Challenge 7
multi-track case with a laser turnaround time of 0.75 ms using the
proposed Double Conical Heat Source (DCHS) model. This case constitutes
the second primary multi-track prediction case of the AM-Bench 2025
competition and was simulated using the DCHS model calibrated from the
corresponding single-track case provided in Online Resource 5. The video
presents the transient evolution of the temperature field, liquid
fraction, volumetric laser heat-source intensity, and radiative and
convective heat-loss source terms during 15 bidirectional tracks. The
simulation was performed with a laser power of 285 W, a scan speed of
960 mm/s, a track length of 1 mm, a hatch spacing of 110~$\mu$m, and a
laser spot diameter of 72~$\mu$m. A turnaround time of 0.75 ms was
applied between consecutive tracks. Detailed computational and processing
parameters are provided on the opening screen of the video.

\medskip

\section*{S2. Temperature-History Videos}

Online Resources 8--14 present the transient temperature histories
extracted from the corresponding single-track and multi-track simulations.
These videos illustrate the heating and cooling cycles experienced at the
selected monitoring locations as the laser passes over or near each point.
The temperature histories are used to evaluate the thermal metrics reported
in the main manuscript, including the Solid Cooling Rate (SCR), Transition
Cooling Rate (TCR), Liquid Cooling Rate (LCR), and Time Above Melt (TAM)
for the single-track cases, as well as the Pad Solid Cooling Rate (PSCR)
and Pad Time Above Melt (PTAM) for the multi-track cases.

The horizontal threshold lines shown in the temperature-history plots
indicate the temperature ranges used to evaluate the corresponding thermal
metrics. Temperature is reported in K and time in seconds. The animations
show the evolution of the temperature history as the transient simulation
progresses and provide a visual representation of the thermal cycles used
for the quantitative analyses presented in the main manuscript.

\textbf{Online Resource 8 (ESM\_8.mp4).}\\
Temperature-history animation for the NIST AM-Bench 2022 single-track
Case 0 simulated using the conventional single-conical heat-source model.
The temperature history corresponds to the single-track validation analysis
presented in the main manuscript and is used to evaluate the associated
cooling-rate and time-above-melt metrics.

\medskip

\textbf{Online Resource 9 (ESM\_9.mp4).}\\
Temperature-history animation for the NIST AM-Bench 2022 single-track
Case 0 simulated using the proposed Double Conical Heat Source (DCHS)
model. The temperature history corresponds to the DCHS single-track
validation presented in the main manuscript and is used to evaluate
the associated Solid Cooling Rate (SCR), Transition Cooling Rate (TCR),
Liquid Cooling Rate (LCR), and Time Above Melt (TAM). This resource
provides the thermal-history counterpart to the transient simulation
presented in Online Resource 2.

\medskip

\textbf{Online Resource 10 (ESM\_10.mp4).}\\
Temperature-history animation for the NIST AM-Bench 2022 XPad
multi-track benchmark simulated using the conventional single-conical
heat-source model. The temperature histories correspond to the XPad
multi-track validation presented in the main manuscript and illustrate
the repeated thermal cycles produced by successive laser tracks. The
histories are used to determine the Pad Time Above Melt (PTAM) and Pad
Solid Cooling Rate (PSCR). This resource provides the temperature-history
counterpart to the transient simulation presented in Online Resource 3.

\medskip

\textbf{Online Resource 11 (ESM\_11.mp4).}\\
Temperature-history animation for the NIST AM-Bench 2022 XPad
multi-track benchmark simulated using the proposed Double Conical
Heat Source (DCHS) model. The temperature histories correspond to
the DCHS XPad multi-track validation presented in the main manuscript
and illustrate the repeated heating and cooling cycles caused by
successive laser tracks and remelting. These histories are used to
determine the Pad Time Above Melt (PTAM) and Pad Solid Cooling Rate
(PSCR). This resource provides the temperature-history counterpart
to the transient DCHS simulation presented in Online Resource 4.

\medskip

\textbf{Online Resource 12 (ESM\_12.mp4).}\\
Temperature-history animation for the NIST AM-Bench 2025 single-track
case simulated using the proposed Double Conical Heat Source (DCHS)
model. This single-track case was used to calibrate the DCHS model
prior to its application to the AM-Bench 2025 Challenge 7 multi-track
prediction cases presented in the main manuscript. The animation
illustrates the transient heating and cooling response at the selected
monitoring location and the temperature thresholds used to evaluate
the corresponding thermal metrics. This resource provides the
temperature-history counterpart to the transient simulation presented
in Online Resource 5.

\medskip

\textbf{Online Resource 13 (ESM\_13.mp4).}\\
Temperature-history animation for the NIST AM-Bench 2025 Challenge 7
multi-track case with a laser turnaround time of 5 ms, simulated using
the proposed Double Conical Heat Source (DCHS) model. This case is one
of the primary multi-track prediction cases of the AM-Bench 2025
competition. The animation illustrates the repeated heating, remelting,
and cooling cycles at the selected monitoring locations during
successive laser tracks. The temperature histories are used to evaluate
the Pad Time Above Melt (PTAM) and Pad Solid Cooling Rate (PSCR)
reported in the main manuscript. This resource provides the
temperature-history counterpart to the transient simulation presented
in Online Resource 6.

\medskip

\textbf{Online Resource 14 (ESM\_14.mp4).}\\
Temperature-history animation for the NIST AM-Bench 2025 Challenge 7
multi-track case with a laser turnaround time of 0.75 ms, simulated
using the proposed Double Conical Heat Source (DCHS) model. This case
is the second primary multi-track prediction case of the AM-Bench 2025
competition. The animation illustrates the repeated heating, remelting,
and cooling cycles at the selected monitoring locations during
successive laser tracks. The temperature histories are used to evaluate
the Pad Time Above Melt (PTAM) and Pad Solid Cooling Rate (PSCR)
reported in the main manuscript. This resource provides the
temperature-history counterpart to the transient simulation presented
in Online Resource 7.

\medskip

\section*{S3. Supporting Numerical Data}

\textbf{Online Resource 15 (ESM\_15.xlsx).}\\
Spreadsheet containing the simulated melt-pool and pad-dimension data for the
NIST AM-Bench 2022 XPad multi-track case. The file provides track-wise depth,
overlap depth, width, and overlap width predicted using both the conventional
single-conical and double conical heat source models. It also includes
comparisons of the predicted pad geometry and thermal metrics with the
experimental measurements, together with the corresponding percentage errors
and experimental standard deviations. The dataset supports the AM-Bench 2022
XPad multi-track validation and comparison presented in the main manuscript.

\medskip

\textbf{Online Resource 16 (ESM\_16.xlsx).}\\
Spreadsheet containing the experimental and simulated pad-dimension data for
the NIST AM-Bench 2025 Challenge 7 multi-track cases with laser turnaround
times of 0.75 ms and 5 ms. The file provides the pad depth, overlap depth,
width, and overlap width at the measurement locations for comparison between
the experimental measurements and predictions obtained using the Double
Conical Heat Source (DCHS) model. These data support the validation of the
AM-Bench 2025 multi-track predictions presented in the main manuscript.

\medskip

\textbf{Online Resource 17 (ESM\_17.xlsx).}\\
Spreadsheet containing the point-wise thermal metrics obtained from the
NIST AM-Bench 2022 single-track Case 0 simulations using the conventional
single-conical and proposed Double Conical Heat Source (DCHS) models.
Separate worksheets provide the Solid Cooling Rate (SCR), Transition
Cooling Rate (TCR), Liquid Cooling Rate (LCR), and Time Above Melt (TAM)
calculated at the selected monitoring points for each heat-source model.
These data support the single-track thermal validation and comparison of
the two heat-source models presented in the main manuscript.

\medskip

\textbf{Online Resource 18 (ESM\_18.xlsx).}\\
Spreadsheet containing the point-wise thermal metrics obtained from the
NIST AM-Bench 2022 XPad multi-track simulations using the conventional
single-conical and proposed Double Conical Heat Source (DCHS) models.
The file provides the Time Above Melt (TAM) and Solid Cooling Rate (SCR)
calculated at the selected monitoring points for both heat-source models.
These data were used to evaluate the Pad Time Above Melt (PTAM) and Pad
Solid Cooling Rate (PSCR) and support the multi-track thermal validation
and comparison presented in the main manuscript.

\medskip

\textbf{Online Resource 19 (ESM\_19.xlsx).}\\
Spreadsheet containing the point-wise thermal metrics obtained from the
NIST AM-Bench 2025 single-track simulation using the proposed Double
Conical Heat Source (DCHS) model. The file provides the Solid Cooling
Rate (SCR), Transition Cooling Rate (TCR), Liquid Cooling Rate (LCR),
and Time Above Melt (TAM) calculated at the selected monitoring points.
These data characterize the thermal response of the calibrated
single-track model used before its application to the AM-Bench 2025
Challenge 7 multi-track prediction cases.

\medskip

\textbf{Online Resource 20 (ESM\_20.xlsx).}\\
Spreadsheet containing the point-wise thermal metrics obtained from the
NIST AM-Bench 2025 Challenge 7 multi-track simulations with laser turnaround
times of 5 ms and 0.75 ms using the proposed Double Conical Heat Source
(DCHS) model. Separate worksheets provide the monitoring-point coordinates,
Time Above Melt (TAM), and Solid Cooling Rate (SCR) for the two turnaround-time
cases. These data were used to determine the Primary Time Above Melt (PTAM)
and Primary Solid Cooling Rate (PSCR) and to evaluate the influence of laser
turnaround time on the multi-track thermal response presented in the main
manuscript.

\medskip

\textbf{Online Resource 21 (ESM\_21.xlsx).}\\
Spreadsheet containing the grid-independence study for the NIST AM-Bench
2022 single-track Case 0 simulation using the proposed Double Conical Heat
Source (DCHS) model. The file compares four computational meshes with
different grid resolutions and cell counts. Melt-pool depth, width, area,
maximum temperature, Solid Cooling Rate (SCR), Transition Cooling Rate
(TCR), Liquid Cooling Rate (LCR), and Time Above Melt (TAM) are reported
for each mesh. These results were used to evaluate mesh sensitivity and
select an appropriate grid resolution for the numerical simulations
presented in the main manuscript.

\medskip

\textbf{Online Resource 22 (ESM\_22.xlsx).}\\
Spreadsheet containing the coordinates of the monitoring points used to
extract temperature histories from the numerical simulations. Separate
worksheets provide the monitoring-point coordinates for the NIST AM-Bench
2025 multi-track, AM-Bench 2022 multi-track, and AM-Bench 2022/2025
single-track cases. The extracted temperature histories were used to plot
the cooling curves and calculate the cooling-rate and time-above-melt
thermal metrics reported in the main manuscript. The $x$, $y$, and $z$
coordinates are reported in meters (m).

\section*{S4. Supporting graphs}

\textbf{Online Resource 23 (ESM\_23\_steady condition.png).}\\

Image showing the evolution of the predicted melt-pool dimensions with
laser travel distance. The melt-pool length, width, and depth are plotted
to demonstrate the development of a steady-state melt pool during
single-track laser scanning. After the initial transient region, the
melt-pool dimensions remain nearly constant with increasing laser travel
distance, indicating that a steady-state condition has been reached. This
analysis supports the selection of the track length used for evaluating
the steady-state melt-pool dimensions in the simulations.

\section*{Summary of Supplementary Files}

\begin{longtable}{p{0.10\textwidth} p{0.18\textwidth} p{0.18\textwidth} p{0.45\textwidth}}
\toprule
\textbf{Resource} & \textbf{File Type} & \textbf{Case} & \textbf{Content} \\
\midrule
\endfirsthead

\toprule
\textbf{Resource} & \textbf{File Type} & \textbf{Case} & \textbf{Content} \\
\midrule
\endhead

ESM\_1 &
Video (MP4) &
AM-Bench 2022, single track &
Transient simulation using the single-conical heat-source model. \\

ESM\_2 &
Video (MP4) &
AM-Bench 2022, single track &
Transient simulation using the Double Conical Heat Source (DCHS) model. \\

ESM\_3 &
Video (MP4) &
AM-Bench 2022, XPad &
Transient 10-track simulation using the single-conical heat-source model. \\

ESM\_4 &
Video (MP4) &
AM-Bench 2022, XPad &
Transient 10-track simulation using the DCHS model. \\

ESM\_5 &
Video (MP4) &
AM-Bench 2025, single track &
Transient DCHS simulation used for single-track model calibration. \\

ESM\_6 &
Video (MP4) &
AM-Bench 2025, Challenge 7 &
Transient 15-track DCHS simulation with 5 ms turnaround time. \\

ESM\_7 &
Video (MP4) &
AM-Bench 2025, Challenge 7 &
Transient 15-track DCHS simulation with 0.75 ms turnaround time. \\

\midrule

ESM\_8 &
Video (MP4) &
AM-Bench 2022, single track &
Temperature history for the single-conical heat-source simulation. \\

ESM\_9 &
Video (MP4) &
AM-Bench 2022, single track &
Temperature history for the DCHS simulation. \\

ESM\_10 &
Video (MP4) &
AM-Bench 2022, XPad &
Temperature histories for the single-conical multi-track simulation. \\

ESM\_11 &
Video (MP4) &
AM-Bench 2022, XPad &
Temperature histories for the DCHS multi-track simulation. \\

ESM\_12 &
Video (MP4) &
AM-Bench 2025, single track &
Temperature history for the DCHS single-track simulation. \\

ESM\_13 &
Video (MP4) &
AM-Bench 2025, Challenge 7 &
Temperature histories for the 5 ms turnaround-time case. \\

ESM\_14 &
Video (MP4) &
AM-Bench 2025, Challenge 7 &
Temperature histories for the 0.75 ms turnaround-time case. \\

\midrule

ESM\_15 &
Spreadsheet (XLSX) &
AM-Bench 2022, XPad &
Pad depth, overlap depth, width, and overlap-width data and comparison. \\

ESM\_16 &
Spreadsheet (XLSX) &
AM-Bench 2025, Challenge 7 &
Experimental and simulated pad dimensions for the 5 ms and 0.75 ms cases. \\

ESM\_17 &
Spreadsheet (XLSX) &
AM-Bench 2022, single track &
SCR, TCR, LCR, and TAM for the single-conical and DCHS models. \\

ESM\_18 &
Spreadsheet (XLSX) &
AM-Bench 2022, XPad &
Point-wise TAM and SCR data for the single-conical and DCHS models. \\

ESM\_19 &
Spreadsheet (XLSX) &
AM-Bench 2025, single track &
Point-wise SCR, TCR, LCR, and TAM from the DCHS simulation. \\

ESM\_20 &
Spreadsheet (XLSX) &
AM-Bench 2025, Challenge 7 &
Point-wise TAM and SCR data for the 5 ms and 0.75 ms turnaround cases. \\

ESM\_21 &
Spreadsheet (XLSX) &
AM-Bench 2022, single track &
Grid-independence results based on geometric and thermal metrics. \\

ESM\_22 &
Spreadsheet (XLSX) &
AM-Bench 2022/2025 &
Coordinates of monitoring points used for temperature-history extraction. \\

\midrule

ESM\_23 &
Image (PNG) &
Single-track simulation &
Melt-pool dimensions versus laser travel distance demonstrating steady state. \\

\bottomrule
\end{longtable}

\section*{Data Availability Note}

The supplementary files provide the numerical and visual material directly
supporting the analyses presented in the manuscript. Additional simulation
files and related data are available through the project repository and may
also be obtained from the corresponding author upon reasonable request.

\end{document}
